\scanpage{19}
% archive foliation: 18
dans la suite des $\operatorname{div}n_{i\alpha}$ est $\beta_{i\ell}$,
et $\Phi(i,\ell)_*$ est le \deleted{complexe}
\inserted{module $\frac12$ simplicial} associé
à
\diagram{\begin{tikzcd}
-0\arrow[r]&\overset{\text{deg $i$}}k\arrow[r,"n"]&
\overset{\text{deg $i-1$}}k\arrow[r]&0-
\end{tikzcd}}
où $\operatorname{div}(n)=\ell$. On a donc
une extension
\diagram{\begin{tikzcd}
0\arrow[r]&\Lambda^{i-1}\Psi\arrow[r]&\Phi(i,n)\arrow[r]&\Lambda^i\Psi\arrow[r]&0
\end{tikzcd}}
où l'hom. assoc. à cette extension
\diagram{\begin{tikzcd}k\arrow[r]&k\end{tikzcd}}
(cf \unsure{\ldots}) est $n\operatorname{id}$, où $\operatorname{div}(n)=\ell$;
donc
\[
\deleted{\unsure{\beta}}\pi_j(\Phi(i,n))=
\left\{\begin{array}{ll}0&\text{si }j\neq i-1\\
k/\operatorname{div}n&\text{si }j=i-1.
\end{array}\right.
\]

\uline{Remarque} 1) La fonction de Hilbert de $\Phi(i,n)$
est donc la même que celle de $\Lambda^i\Phi_*$,
savoir $\binom{n+1}{i}$. Donc si $L_*$ est
donné par le produit de la prop on a
\[
\operatorname{rg}L_n=\sum_{i\leq n+1}(\mu_i+\beta_i)\binom{n+1}{i}
+\sum_{0\leq j\leq n}\nu_j\binom{n}{\unsure{i|j}}.
\]
Donc si $L_*$ est \inserted{$\partial$-parfait, il est}
connu à isom. (non unique)
près quand on connaît a) $\unsure{\operatorname{rg}\pi_i(L_*)=\nu_i}$
et le nombre $\beta_{i+1}$ de ses diviseurs de
torsion, pour tt $i\geq0$, et b) la fonction
$f_{L_*}(n)$ pour $n$ grand (qui donne
\deleted{\unsure{alors}} \inserted{pour $n$ grand} le polynôme
$\sum_{i\geq1}(\mu_i+\beta_i)\binom{n+1}{i}
+\sum_{j\geq0}\nu_j\binom{n}{\unsure{i|j}}$)
